\item Un pequeño disco de radio $r$ y masa $m$ se pega rígidamente a la cara de un segundo disco más grande de radio $R$ y masa $M$, como se muestra en la figura P15.48. El centro del disco pequeño se ubica en el borde del disco grande. El disco grande se monta en su centro en un eje sin fricción. El ensamble da vueltas a través de un pequeño ángulo $\theta$ desde su posición de equilibrio y se libera.
\begin{enumerate}
    \item Demuestre que mientras pasa a través de la posición de equilibrio la rapidez del centro del disco pequeño es:
    $$ v = 2 \left[ \frac{Rg(1 - \cos\theta)}{3(M/m) + (r/R)^2 + 2} \right]^{1/2} $$
    \item Demuestre que el periodo del movimiento es:
    $$ T = 2\pi \left[ \frac{3(M + 2m)R^2 + 2mr^2}{4mgR} \right]^{1/2} $$
\end{enumerate}